\documentclass[10pt,a4paper]{amsart}
\usepackage[margin=19mm]{geometry}
\usepackage{amsmath,amssymb,mathtools}
\usepackage[scr=rsfs,cal=euler]{mathalfa}
\usepackage{xfrac}
\usepackage[unicode,hidelinks,bookmarks=false]{hyperref}
\newcommand{\ie}{i.e.\ }
\newcommand{\cf}{cf.\ }
\newcommand{\resp}{respectively\ }
\newcommand{\Subsection}{\subsection}
\providecommand{\nobreakdash}{\nobreak\hbox{-}}
\setlength{\emergencystretch}{2em}
\multlinegap0pt
\usepackage{bm}
\usepackage{bbm}
\usepackage{dsfont}
\usepackage{xspace}
\usepackage{cancel}
\usepackage{mathtools}
\usepackage{amsthm}
\makeatletter
\@ifundefined{lemma}{\newtheorem{lemma}{Lemma}}{}
\makeatother
\title[Path-Based Conditions for the Identifiability of Non-additiv]{Path-Based Conditions for the Identifiability of Non-additive Nonlinear Networks with Full Measurements}
\author{Renato Vizuete, Julien M. Hendrickx}
\date{Source: arXiv:2510.20537v2 (CC BY 4.0)}
\begin{document}
\maketitle

\noindent\emph{Source and licence.} Path-Based Conditions for the Identifiability of Non-additive Nonlinear Networks with Full Measurements. Version arXiv:2510.20537v2 (CC BY 4.0), distributed under the Creative Commons Attribution 4.0 International licence, as recorded in the arXiv licence metadata for this exact version and shown on the abstract page https://arxiv.org/abs/2510.20537v2. Full rights notice: licenses/arxiv-2510.20537-CC-BY-4.0.txt.\par\smallskip

\noindent\textit{Editorial scope.} Counted unit: the Lemma whose source label is \texttt{lemma:sinks\_sources} in the article (source0.tex lines 468--470, source label \texttt{lemma:sinks\_sources}), delivered together with the article's own complete original proof of it (source0.tex lines 471--486, one \texttt{proof} environment of 16 lines). No mathematics is written, repaired or re-derived: statement and proof are the article's own text line for line. Required context retained uncounted: none. Every definition, display and label of those lines is retained unchanged. Omitted from this package and not redistributed: 1--467, 487--781. Nothing inside a retained excerpt is omitted or altered: every retained body is delivered exactly as the article's lines are written, so no omission receipt of any kind accompanies this package. Citations: the retained excerpts carry no citation command at all, so no bibliography entry of the article is transcribed and no reference is redistributed. Reference adaptation: none was needed and none was made; every reference inside the delivered unit resolves inside it, so no unresolved reference or citation is produced. Printed slips kept literal: no printed word, symbol, hypothesis or number of the counted statement or of its proof was altered. Layout and class: the article's own class is \texttt{IEEEtran} with packages including graphics, epsfig, amsmath, amssymb, dsfont, enumerate, subfigure, multirow, bm, tikz, cite, amsthm, array, makecell; the wrapper is the lane's pinned \texttt{amsart} editorial wrapper at 10pt on A4, which restates the theorem-like environments the retained text needs and adds the article's own macro definitions that the retained text needs (none), with extra wrapper packages (bm, bbm, dsfont, xspace, cancel, mathtools). The delivered numbering is the unit's own. Assets: no retained span contains a figure environment, a graphics-inclusion command, a PDF-page inclusion, a table or a TikZ picture; no image file is copied into this package, the package declares no asset dependency and no image file is redistributed with it. Licence: version arXiv:2510.20537v2 of this article is distributed under the Creative Commons Attribution 4.0 International licence, as recorded in the arXiv licence metadata for this exact version and shown on the abstract page https://arxiv.org/abs/2510.20537v2. A full rights notice is delivered at licenses/arxiv-2510.20537-CC-BY-4.0.txt. Characters: every retained excerpt is pure ASCII, so the delivered engine, XeLaTeX with its default Latin Modern text font, reports no missing character.
\begin{lemma}\label{lemma:sinks_sources}
    The outgoing edges of sources are not identifiable if the sources are not excited. The excitation of sinks is never necessary for the identifiability of the network. The measurement of sources is never necessary.
\end{lemma}

\begin{proof}
Without loss of generality, let us consider that the node 1 is a source and $j$ is an out-neighbor of 1.
    The measurement of $j$ provides the output
    \begin{equation}\label{eq:excitation_source1}
        y_j^k=u_j^{k-1}+\Phi_j(y_1^{k-0},\ldots,y_1^{k-m_{j,1}},y_2^{k-0},\ldots,y_{M_j}^{k-m_{j,M_j}}).    
    \end{equation}
    If the source $1$ is not excited (i.e., $u_1^k=0$ for all $k$), its output $y_1^k=0$, and   \eqref{eq:excitation_source1} becomes
    \begin{equation}\label{eq:excitation_source2}
        y_j^k=u_j^{k-1}+\Phi_j(0,\ldots,0,y_2^{k-0},\ldots,y_{M_j}^{k-m_{j,M_j}}).  
    \end{equation}
    Notice that any node function $\tilde \Phi_j=\psi_1(y_1^{k-0},\ldots,y_1^{k-m_{j,1}})+\Phi_j$ where $\psi_1$ is analytic and  $\psi_1(0)=0$, also satisfies \eqref{eq:excitation_source2}. This implies that it is not possible to identify $\Phi_j$ and hence, the excitation of all the sources is necessary for identifiability of the network.
    Now, the excitation of a sink $\ell$ can only affect its output:
    $y_\ell^k=u_\ell^{k-1}+\Phi_\ell$. Therefore,
    the identifiability of $\Phi_\ell$ is independent of the value of $u_\ell^{k-1}$, which implies that the excitation of the sink is never necessary for identifiability.
Finally, the measurement of a source $i$ provides the output $y_i=u_i^{k-1}$ that does not include any unknown dynamics.
\end{proof}

\end{document}
